% SYNTHETIC EVALUATION DRAFT — written for the retrieval eval in evals/.
% Not a real paper; contains no experimental results.
\documentclass{article}
\title{Rank-Adaptive Low-Rank Adapters for Parameter-Efficient Fine-Tuning of Large Language Models}
\begin{document}
\maketitle

\begin{abstract}
Full fine-tuning of large language models updates billions of parameters and
requires storing a separate copy of the model for every downstream task.
Parameter-efficient fine-tuning (PEFT) methods such as LoRA, adapters, prefix
tuning and BitFit freeze the pretrained backbone and train only a small number
of additional parameters. This draft proposes a rank-adaptive variant of
low-rank adaptation that allocates the rank budget across transformer layers
according to a sensitivity score computed during training, and studies how it
interacts with 4-bit quantized backbones.
\end{abstract}

\section{Introduction}
Pretrained transformer language models are routinely adapted to instruction
following, summarization, code generation and domain-specific question
answering. Updating all weights of a model with tens of billions of parameters
is expensive in GPU memory, optimizer state and storage. Parameter-efficient
fine-tuning addresses this by keeping the pretrained weights frozen and
learning a compact task-specific update.

Low-rank adaptation (LoRA) injects a trainable low-rank decomposition
$\Delta W = BA$ into the attention projection matrices, where
$B \in \mathbb{R}^{d \times r}$, $A \in \mathbb{R}^{r \times k}$ and the rank
$r$ is much smaller than $d$ and $k$. After training the update can be merged
into the frozen weight, so inference latency is unchanged. A known limitation
is that a single global rank is used for every layer, although layers differ
in how much they need to change for a given task.

Our contributions are: (1) a sensitivity-guided rank allocation scheme that
grows and prunes the rank of each LoRA module during fine-tuning under a fixed
total parameter budget; (2) an analysis of rank allocation when the backbone is
quantized to 4-bit NormalFloat, as in QLoRA; and (3) a planned evaluation
protocol on instruction tuning and commonsense reasoning benchmarks.

\section{Related Work}
Adapters insert small bottleneck feed-forward modules between transformer
sublayers. Prefix tuning and prompt tuning prepend trainable continuous vectors
to the keys and values or to the input embeddings. BitFit trains only the bias
terms. LoRA and its variants reparameterize the weight update as a low-rank
product. AdaLoRA parameterizes the update with a singular value decomposition
and prunes singular values by importance. QLoRA backpropagates through a frozen
4-bit quantized model into LoRA adapters, enabling fine-tuning of very large
models on a single GPU. DoRA decomposes weights into magnitude and direction
components. Our work is closest to AdaLoRA but differs in that ranks can also
grow, not only shrink.

\section{Method}
\subsection{Sensitivity Score}
For each LoRA module $m$ we maintain an exponential moving average of
$|\theta \cdot \nabla_\theta \mathcal{L}|$ over its parameters, which
approximates the change in loss if the module were removed. Modules with high
sensitivity are candidates for rank growth.

\subsection{Rank Budget Reallocation}
Every $T$ optimizer steps, we rank all modules by sensitivity. The lowest
scoring modules have their rank reduced by pruning the column of $B$ and the
row of $A$ with the smallest norm product; the freed budget is given to the
highest scoring modules by appending newly initialized rank-one components,
with $B$ initialized to zero so that the model output is unchanged at the moment
of growth. The total number of trainable parameters is held constant.

\subsection{Interaction with Quantization}
When the frozen backbone is quantized, quantization error acts as a fixed
perturbation of the pretrained weights. We hypothesize that layers with larger
quantization error receive larger rank under our scheme, and we plan to
measure the correlation between per-layer quantization error and allocated
rank.

\section{Experimental Plan}
We plan to fine-tune decoder-only models of several sizes on instruction
tuning data and evaluate on commonsense reasoning and reading comprehension
benchmarks. Baselines are full fine-tuning, LoRA with a uniform rank, AdaLoRA,
adapters and prefix tuning, all at matched trainable parameter counts. We will
report task accuracy, peak GPU memory, number of trainable parameters and
wall-clock training time. Ablations will vary the reallocation interval $T$,
the total rank budget and whether growth is enabled.

\section{Limitations}
The sensitivity score adds a small overhead per step. Rank changes alter the
optimizer state shape, which requires resetting the Adam moments of the
affected components. The approach has only been designed for attention and
MLP projection matrices in transformer language models.

\end{document}
